% !TeX root = ../../../main.tex
\subsection{Validity of Temporal Annotations}
\label{ch07:apd:tts_validity}


We used \texttt{DeepSeek-R1-UD-IQ1} and \texttt{Llama-3.3-70B-Instruct} models to generate clinical textual time series from
the 2319 sepsis case reports.
We refer to this as Textual
Time Series for Sepsis (\textit{T2S2}).
We assessed the quality of \textit{T2S2} using three key metrics: event match rate, temporal concordance (c-index), and log-time discrepancy on two external validation sets: \textit{sepsis-10} (n = 10) and \textit{sepsis-100} (n = 90).(adopted from our prior work \cite{wang2025large}).

\ChSevenStarSubsection{Text Event Match}
To align predicted clinical findings with reference annotations, we use a recursive best-match procedure adapted from \citep{wang2025large}.
At each step, we select the reference-prediction pair with the minimal string-based distance; ties are broken by text order (\emph{i.e.}, the order of events in the annotation files).
Matched items are removed and the process recurses on the remainder, preventing many-to-one mappings.
We evaluated Levenshtein distance, cosine similarity of \textit{BERT} (\textit{bert-base-uncased}) embeddings, and sentence-embedding cosine similarity from \textit{PubMedBERT} (\textit{S-PubMedBert-MS-MARCO}).
Prior validation showed \textit{PubMedBERT-based} cosine similarity with a threshold of 0.1 performed best, and we used this metric for evaluation.
Pseudocode is provided in Algorithm \ref{ch07:alg:recursive_match}.


\ChSevenStarSubsection{Temporal Ordering}
Temporal ordering accuracy was quantified with the concordance index (c-index), the probability that a randomly chosen pair of events is correctly ordered.
For reference times $\{t^r_i\}$ and predicted times $\{t^p_i\}$,
\[
\text{c-index}=\frac{1}{N}\sum_{i<j}\mathds{1}\{(t^r_i-t^r_j)(t^p_i-t^p_j)>0\},
\]
where $N$ is the number of comparable pairs (event indices $i,j$ with $t^r_i\neq t^r_j$ and $t^p_i\neq t^p_j$), and $\mathds{1}\{\cdot\}$ is the indicator function.
Higher values indicate better agreement with the reference timeline.
 



% !TeX root = ../../../main.tex
\begin{ChSevenAlgorithm}[tbp]
\caption{Recursive Best Match}
\label{ch07:alg:recursive_match}
\begin{algorithmic}[1]
\ChSevenAlgRequire Two lists: \texttt{ref} (reference events) and \texttt{pred} (predicted events)
\ChSevenAlgEnsure List of best-matching event pairs
\ChSevenAlgFunction{MatchEvents}{\texttt{ref}, \texttt{pred}}
    \ChSevenAlgIf{\texttt{ref} is empty \textbf{or} \texttt{pred} is empty}
        \ChSevenAlgState \ChSevenAlgReturn $[]$
    \ChSevenAlgEndIf
    \ChSevenAlgState $\text{min\_distance} \gets \infty$
    \ChSevenAlgState $\text{best\_pair} \gets \text{None}$
    \ChSevenAlgForAll{$r \in \texttt{ref}$}
        \ChSevenAlgForAll{$p \in \texttt{pred}$}
            \ChSevenAlgState $d \gets \text{ComputeDistance}(r, p)$
            \ChSevenAlgIf{$d < \text{min\_distance}$}
                \ChSevenAlgState $\text{min\_distance} \gets d$
                \ChSevenAlgState $\text{best\_pair} \gets (r, p)$
            \ChSevenAlgElsIf{$d = \text{min\_distance}$}
                \ChSevenAlgState Get indices of $r$, $p$, and current $\text{best\_pair}$
                \ChSevenAlgIf{index of $r <$ best $r$ index}
                    \ChSevenAlgState $\text{best\_pair} \gets (r, p)$
                \ChSevenAlgElsIf{index of $r = $ best $r$ index \textbf{and} index of $p <$ best $p$ index}
                    \ChSevenAlgState $\text{best\_pair} \gets (r, p)$
                \ChSevenAlgEndIf
            \ChSevenAlgEndIf
        \ChSevenAlgEndFor
    \ChSevenAlgEndFor
    \ChSevenAlgState Remove $\text{best\_pair.r}$ from \texttt{ref}
    \ChSevenAlgState Remove $\text{best\_pair.p}$ from \texttt{pred}
    \ChSevenAlgState $\text{result} \gets [\text{best\_pair}] + \text{MatchEvents}(\texttt{ref}, \texttt{pred})$
    \ChSevenAlgState \ChSevenAlgReturn $\text{result}$
\ChSevenAlgEndFunction
\end{algorithmic}
\end{ChSevenAlgorithm}



\ChSevenStarSubsection{Time Discrepancy Analysis}
 Temporal prediction accuracy was assessed via discrepancies between predicted and reference timestamps.
We computed the log-time discrepancy as $\log(1+\Delta t)$, where $\Delta t$ is the absolute time difference (in hours), and summarized errors with the Area Under the Log-Time Cumulative Distribution Function (AULTC), which reflects how rapidly the error CDF rises—larger AULTC indicates better temporal alignment. 


% !TeX root = ../../../main.tex
\begin{table}[htbp]
\caption{Comparison of event match rate (Event), median concordance (c), median absolute error (MAE), and area under the log-time curve (AULTC) across \texttt{DeepSeek-R1} and \texttt{Llama-3.3-70B} for \textit{sepsis-10}, and \textit{sepsis-100} datasets.}\ChSevenFloatLabel{ch07:tab:sepsis_onecol}
\ChSevenTableBodyStart
\begin{ChSevenFloatBody}

\centering
\begin{adjustbox}{max width=\columnwidth}
\setlength{\tabcolsep}{4.5pt} %
\renewcommand{\arraystretch}{1.1}

\begin{tabular}{lccc}
\hline
\multicolumn{4}{c}{\textbf{\textit{sepsis-10}}} \\
\hline
\textbf{Model} & \textbf{Event Match Rate} & \textbf{Concordance} & \textbf{AULTC} \\
\texttt{DeepSeek-R1}      & 0.79 & 0.876 & 0.772 \\
\texttt{Llama-3.3-70B}    & 0.75 & 0.932 & 0.759 \\
\hline
\multicolumn{4}{c}{\textbf{\textit{sepsis-100}}} \\
\hline
\textbf{Model} & \textbf{Event Match Rate} & \textbf{Concordance} & \textbf{AULTC} \\
\texttt{DeepSeek-R1}      & 0.91 & 0.875 & 0.801 \\
\texttt{Llama-3.3-70B}    & 0.78 & 0.894 & 0.791 \\
\hline
\end{tabular}
\end{adjustbox}



\end{ChSevenFloatBody}
\end{table}



Formally, we compute the AULTC as the area under $F(x)$ from $x = 0$ to $x = \log(1 + S_{\text{max}})$, normalized by $\log(1 + S_{\text{max}})$:

\begin{equation*}
\ChSevenUnusedEquationLabel{ch07:eq:aultc}
\begin{aligned}
\text{AULTC}
&= \frac{1}{\log(1+S_{\max})}
   \sum_{i=1}^{k}\big(x_{(i)}-x_{(i-1)}\big)\frac{i}{k} \\
&\quad + \frac{\log(1+S_{\max})-x_{(k)}}{\log(1+S_{\max})},
\end{aligned}
\end{equation*}

where $x_{(0)} = 0$.
With this definition, AULTC = 1 indicates that discrepancies are zero (perfect recovery), resulting in maximum area $\log(1 + S_{\text{max}})$, and AULTC = 0 indicates that all discrepancies exceed $S_{\text{max}}$, yielding zero area.
We also stratified results by reference time intervals (hour, day, week, year) to examine accuracy as a function of time from case presentation.

\ChSevenStarSubsection{Annotation Quality Results}
On the \textit{sepsis-10} subset, \texttt{DeepSeek-R1} achieved an event match rate of 0.79, c-index of 0.876, and AULTC of 0.772, while \texttt{Llama-3.3-70B} obtained 0.75, 0.932, and 0.759, respectively.
Thus, \texttt{Llama} provided stronger pairwise ordering (higher c-index) but with fewer matched events and a slightly weaker log-time error profile (lower AULTC) than \texttt{DeepSeek}.
On \textit{sepsis-100}, \texttt{DeepSeek-R1} reached 0.91 (event match rate), 0.875 (c-index), and 0.801 (AULTC), whereas \texttt{Llama-3.3-70B} reached 0.78, 0.894, and 0.791.
Here, \texttt{DeepSeek} matched substantially more events and achieved better AULTC, while \texttt{Llama} again showed a small c-index advantage.
Overall, \texttt{DeepSeek} tends to optimize event identification and log-time alignment, whereas \texttt{Llama} slightly favors pairwise temporal ordering. 

Overall, these metrics together indicate strong alignment of the generated textual time series with expert annotations, sufficient to support forecasting benchmarks. 




